% !PRINTED 298
\weberpage{298}
%<<CONT
\[
\begin{pmatrix} 1, & 1 \\ 2, & -2 \end{pmatrix}\begin{pmatrix} 0, & 1 \\ -1, & 0 \end{pmatrix}\begin{pmatrix} -2, & -1 \\ -2, & 1 \end{pmatrix}\begin{pmatrix} 0, & 1 \\ -1, & 0 \end{pmatrix} = \begin{pmatrix} -2, & 0 \\ 0, & 2 \end{pmatrix},
\]
\[
\begin{pmatrix} -2, & 0 \\ 0, & 2 \end{pmatrix}\begin{pmatrix} 1, & 0 \\ -1, & 1 \end{pmatrix}\begin{pmatrix} 2, & 0 \\ 0, & -2 \end{pmatrix}\begin{pmatrix} 1, & 0 \\ 1, & 1 \end{pmatrix} = \begin{pmatrix} 1, & 0 \\ 2, & 1 \end{pmatrix},
\]
d.\,h. $A^2$ und mithin alle Potenzen von $A$, woraus wie oben zu schließen, daß $\mathfrak{N}$ mit $\mathfrak{L}_0$ identisch ist. Wir haben also den Satz:

\gesperrt{Ist $p > 3$, so hat die Gruppe $\mathfrak{L}_0$ keinen Normalteiler.}

Im Falle $p = 3$ enthält $\mathfrak{N}$ gleichfalls
\[
\begin{pmatrix} 0, & 1 \\ -1, & 0 \end{pmatrix},
\]
also auch
\[
\begin{pmatrix} 1, & 0 \\ \pm1, & 1 \end{pmatrix}\begin{pmatrix} 0, & 1 \\ -1, & 0 \end{pmatrix}\begin{pmatrix} 1, & 0 \\ \mp1, & 1 \end{pmatrix} = \begin{pmatrix} \mp1, & 1 \\ 1, & \pm1 \end{pmatrix},
\]
und es bilden auch in der Tat die vier Substitutionen
\weq{15}{\begin{pmatrix} 1, & 0 \\ 0, & 1 \end{pmatrix},\ \begin{pmatrix} 0, & 1 \\ -1, & 0 \end{pmatrix},\ \begin{pmatrix} -1, & 1 \\ 1, & 1 \end{pmatrix},\ \begin{pmatrix} 1, & 1 \\ 1, & -1 \end{pmatrix}}
einen Normalteiler von $\mathfrak{L}_0$ im Index 3.

Im Falle $p = 3$ ist die Gruppe $\mathfrak{L}_0$ isomorph mit der alternierenden Gruppe der Vertauschungen von vier Elementen. Es ist dies die Gruppe einer \gesperrt{beliebigen} Gleichung vierten Grades, wenn die Quadratwurzel aus der Diskriminante dem Rationalitätsbereich adjungiert wird. Der Divisor (15) dieser Gruppe vom Index 3 liefert dann die in der Algebra bekannte kubische Resolvente der biquadratischen Gleichung.

Eine zur Gruppe (15) \gesperrt{gehörige} Funktion der Wurzeln $v_\infty, v_0, v_1, v_2$ der Transformationsgleichung für den dritten Transformationsgrad ist z.\,B.
\[
v_\infty v_0 + v_1 v_2
\]
oder
\[
(v_\infty - v_0)(v_1 - v_2).
\]
Der Unterschied zwischen diesen beiden Funktionen ist der, daß die erstere durch die Substitution der Determinante $-1$
\[
\begin{pmatrix} -1, & 0 \\ 0, & 1 \end{pmatrix}
\]
ungeändert bleibt, während die zweite dabei ihr Zeichen ändert. Die erstere wird daher zu einer Gleichung dritten Grades mit \gesperrt{rationalen} Koeffizienten führen, während in der kubischen Gleichung für die zweite noch $\sqrt{-3}$ auftritt (§ 78).
